\documentclass[10pt,a4paper]{amsart}
\usepackage[margin=19mm]{geometry}
\usepackage{amsmath,amssymb,mathtools}
\usepackage[scr=rsfs,cal=euler]{mathalfa}
\usepackage{xfrac}
\usepackage[unicode,hidelinks,bookmarks=false]{hyperref}
\newcommand{\ie}{i.e.\ }
\newcommand{\cf}{cf.\ }
\newcommand{\resp}{respectively\ }
\newcommand{\Subsection}{\subsection}
\providecommand{\nobreakdash}{\nobreak\hbox{-}}
\setlength{\emergencystretch}{2em}
\multlinegap0pt
\usepackage{bm}
\usepackage{bbm}
\usepackage{dsfont}
\usepackage{xspace}
\usepackage{cancel}
\usepackage{mathtools}
\usepackage{amsthm}
\makeatletter
\@ifundefined{theorem}{\newtheorem{theorem}{Theorem}}{}
\@ifundefined{proposition}{\newtheorem{proposition}[theorem]{Proposition}}{}
\makeatother
\newcommand{\C}{\mathbb C}
\newcommand{\LL}{\mathcal{L}}
\newcommand{\R}{\mathbb R}
\newcommand{\eps}{\varepsilon}
\title[Dynamical Zeta functions for differentiable parabolic maps o]{Dynamical Zeta functions for differentiable parabolic maps of the interval}
\author{Claudio Bonanno, Roberto Castorrini}
\date{Source: arXiv:2403.17700v3 (CC BY 4.0)}
\begin{document}
\maketitle

\noindent\emph{Source and licence.} Dynamical Zeta functions for differentiable parabolic maps of the interval. Version arXiv:2403.17700v3 (CC BY 4.0), distributed under the Creative Commons Attribution 4.0 International licence, as recorded in the arXiv licence metadata for this exact version and shown on the abstract page https://arxiv.org/abs/2403.17700v3. Full rights notice: licenses/arxiv-2403.17700-CC-BY-4.0.txt.\par\smallskip

\noindent\textit{Editorial scope.} Counted unit: the Proposition whose source label is \texttt{prop:l0-spec} in the article (source0.tex lines 396--398, source label \texttt{prop:l0-spec}), delivered together with the article's own complete original proof of it (source0.tex lines 400--445, one \texttt{proof} environment of 46 lines). No mathematics is written, repaired or re-derived: statement and proof are the article's own text line for line. Required context retained uncounted: asintotica-al (lines 167--169). Every definition, display and label of those lines is retained unchanged. Omitted from this package and not redistributed: 1--166, 170--395, 399--399, 446--961. Nothing inside a retained excerpt is omitted or altered: every retained body is delivered exactly as the article's lines are written, so no omission receipt of any kind accompanies this package. Citations: the retained excerpts carry no citation command at all, so no bibliography entry of the article is transcribed and no reference is redistributed. Reference adaptation: none was needed and none was made; every reference inside the delivered unit resolves inside it, so no unresolved reference or citation is produced. Printed slips kept literal: no printed word, symbol, hypothesis or number of the counted statement or of its proof was altered. Layout and class: the article's own class is \texttt{amsart} with packages including geometry, graphicx, amsmath, amssymb, amsthm, enumerate, xcolor, hyperref; the wrapper is the lane's pinned \texttt{amsart} editorial wrapper at 10pt on A4, which restates the theorem-like environments the retained text needs and adds the article's own macro definitions that the retained text needs (\texttt{\textbackslash C}, \texttt{\textbackslash LL}, \texttt{\textbackslash R}, \texttt{\textbackslash eps}), with extra wrapper packages (bm, bbm, dsfont, xspace, cancel, mathtools). The delivered numbering is the unit's own. Assets: no retained span contains a figure environment, a graphics-inclusion command, a PDF-page inclusion, a table or a TikZ picture; no image file is copied into this package, the package declares no asset dependency and no image file is redistributed with it. Licence: version arXiv:2403.17700v3 of this article is distributed under the Creative Commons Attribution 4.0 International licence, as recorded in the arXiv licence metadata for this exact version and shown on the abstract page https://arxiv.org/abs/2403.17700v3. A full rights notice is delivered at licenses/arxiv-2403.17700-CC-BY-4.0.txt. Characters: every retained excerpt is pure ASCII, so the delivered engine, XeLaTeX with its default Latin Modern text font, reports no missing character.
\begin{equation}\label{asintotica-al}
a_\ell \sim \, const(T)\ \ell^{-1/\alpha}\, (1+O(\ell^{-1}))\quad \text{as $\ell \to \infty$.}
\end{equation}

\begin{proposition}\label{prop:l0-spec}
Let $T:[0,1]\to [0,1]$ be a map satisfying (H1)-(H5) with $\alpha\ge 1$, and let $\LL_0$ be the part of the Perron-Frobenius operator of $T$ associated to the parabolic branch, that is $\LL_0 f = \psi'_0 \, (f\circ \psi_0)$. The essential spectral radius of $\LL_0$ on $C^1$ is greater than or equal to 1.
\end{proposition}

\begin{proof}
We show that every $\lambda \in \C$ with $|\lambda|<1$ is an eigenvalue for $\LL_0$ on $C^1$. Let $f\in C^1([0,1])$ and let $F\in C^{2}$ satisfy $F'=f$. Notice that $T$ is $C^2$ at 0 and the same is true for $\psi_0$. Then
\[
\LL_0 f = \lambda \, f \quad \Leftrightarrow \quad \psi'_0 \, (f\circ \psi_0) = \lambda \, f \quad \Leftrightarrow \quad (F\circ \psi_0)' = \lambda\, F'.
\]
Therefore $\lambda \in \C$ is an eigenvalue for $\LL_0$ on $C^1$ if and only if there exists $F\in C^2$ such that
\begin{equation} \label{new-form-l0}
(F\circ \psi_0)(x) = \lambda\, F(x) + const(F),\qquad \forall\, x\in [0,1].
\end{equation}
We show the existence of a solution to \eqref{new-form-l0} for all $\lambda \in \C$ with $|\lambda|<1$.

First of all, since $\psi_0(0)=0$ and $\psi_0'(0)=1$, then a solution $F$ to \eqref{new-form-l0} with $\lambda\not=1$ satisfies $F'(0)=0$, and the same must be true for the second derivative, that is $F''(0)=0$. 

Let now $h:\R \to \R$ be a $C^2$ function with support in $\overline{A_0} =[a,1]$ and such that $h^{(j)}(a) = h^{(j)}(1) =0$ for $j=1,2$. Then for all $\lambda \in \C$ with $|\lambda|<1$
\begin{equation}\label{f-lambda}
F_\lambda(x) := h(x) + \sum_{\ell=1}^{+\infty}\, \lambda^\ell\, \chi_{_{\overline{A_\ell}}}(x)\, h(T^\ell(x)),\qquad \text{for $x\in (0,1]$}
\end{equation}
with $F_\lambda(0)=0$ is a solution to \eqref{new-form-l0} with $const(F)=0$, where $\chi_{A}$ is the indicator function of the set $A$. Since $|\lambda|<1$ we show that the series is absolutely convergent in $C^2$. The function $\chi_{_{\overline{A_\ell}}}(x)\, h(T^\ell(x))$ is in $C^2([0,1])$ for all $\ell\ge 1$, and we can compute
\[
\| \chi_{_{\overline{A_\ell}}}(x)\, h(T^\ell(x)) \|_\infty \le \|h\|_2, \quad \forall\, \ell\ge 1,
\]  
\[
\| \Big( \chi_{_{\overline{A_\ell}}}(x)\, h(T^\ell(x)) \Big)' \|_\infty \le \| h\|_2\, \prod_{i=1}^{\ell}\, \| T' \|_{\infty,A_i},
\]
where $\|\cdot\|_{\infty,A_i}$ denotes the sup-norm of a function restricted to $A_i$, and
\[
\| \Big( \chi_{_{\overline{A_\ell}}}(x)\, h(T^\ell(x)) \Big)'' \|_\infty \le \| h\|_2\, \prod_{i=1}^{\ell}\, \| T' \|_{\infty,A_i}^2\, \Big(1+ \ell \|T\|_2\Big).
\]
By (H5), for $\ell$ big enough so that $a_{\ell-1}<\eps$, we have 
\[
\| T' \|_{\infty,A_\ell} = T'(a_{\ell-1}) \sim 1 + c a_{\ell-1}^\alpha \sim 1 + const(T) \, \ell^{-1} \quad \text{as $\ell \to \infty$,}
\]
using (H4) and \eqref{asintotica-al}, where $const(T)>0$. Hence
\[
\prod_{i=1}^{\ell}\, \| T' \|_{\infty,A_i} \sim \ell^{const(T)} \quad \text{as $\ell \to \infty$.}
\]
This implies that $F_\lambda \in C^2([0,1])$. In addition, using that $\psi_0([0,1]) = [0,a]$, $\psi_0(\overline{A_{\ell-1}}) = \overline{A_\ell}$ and $T^\ell \circ \psi_0(x) = T^{\ell-1}(x)$ for $x\in\overline{A_{\ell-1}}$, we can compute 
\[
\begin{aligned}
F_\lambda(\psi_0(x)) =\, & h(\psi_0(x)) + \sum_{\ell=1}^{+\infty}\, \lambda^\ell\, \chi_{_{\overline{A_\ell}}}(\psi_0(x))\, h(T^\ell(\psi_0(x))) \\
=\, & \sum_{\ell=1}^{+\infty}\, \lambda^\ell\, \chi_{_{\overline{A_{\ell-1}}}}(x)\, h(T^{\ell-1}(x)) \\
=\, & \lambda\, h(x) + \lambda\, \sum_{\ell=1}^{+\infty}\, \lambda^\ell\, \chi_{_{\overline{A_\ell}}}(\psi_0(x)) = \lambda\, F_\lambda(x)
\end{aligned}
\]
for all $x\in (0,1]$. This completes the proof.
\end{proof}

\end{document}
